%----------------------------------------------------------------------------------------
% TRANSFORMATIES VAN GRAFIEKEN. Alle opgaven zijn eigen werk.
%
% Opbouw: Start, leerstof (theorie, werkwijze of uitleg, voorbeeld, korte oefening), en op een
% nieuwe pagina de Opdrachten (door elkaar, stijgende moeilijkheid), Koppel en Test jezelf.
% Labels van dit hoofdstuk beginnen met h2-.
%----------------------------------------------------------------------------------------

\mychapter{2}{Transformaties van grafieken}

\lkr{Leerplan B+S': LPD 5-6 (algemene functies), het deel over transformaties van grafieken. In de cursussen over exponentiële, goniometrische en rationale functies komen dezelfde transformaties terug: daar volstaat dan een verwijzing naar dit hoofdstuk. Laat leerlingen telkens één bijzonder punt volgen (een top, een beginpunt): dat maakt de transformatie zichtbaar en controleerbaar.}

Veel functies zijn een bekende functie in een andere vorm: verschoven, uitgerekt of gespiegeld. Wie de grafieken van een handvol basisfuncties kent, kan zo heel wat grafieken schetsen zonder tabel.

\begin{oefeningen}[Start]
\oef{h2-start-1}

\oef{h2-verken}
\end{oefeningen}

%========================================================================================
\section{Basisfuncties}

Deze grafieken ken je best uit het hoofd, met hun domein en bereik.

\figuren{%
\begin{tikzpicture}[x=.36cm,y=.36cm,font=\scriptsize]
  \roosterleeg{-3}{1}{3}{-3}{1}{3}
  \kromme{-1.73}{1.73}{(\x)^2}
  \node[below,text=black] at (0,-3.2) {$x^2$};
\end{tikzpicture}\hspace{3mm}
\begin{tikzpicture}[x=.36cm,y=.36cm,font=\scriptsize]
  \roosterleeg{-3}{1}{3}{-3}{1}{3}
  \kromme{-1.44}{1.44}{(\x)^3}
  \node[below,text=black] at (0,-3.2) {$x^3$};
\end{tikzpicture}\hspace{3mm}
\begin{tikzpicture}[x=.36cm,y=.36cm,font=\scriptsize]
  \roosterleeg{-3}{1}{3}{-3}{1}{3}
  \kromme{0}{3}{sqrt(\x)}
  \vol{(0,0)}
  \node[below,text=black] at (0,-3.2) {$\sqrt{x}$};
\end{tikzpicture}\hspace{3mm}
\begin{tikzpicture}[x=.36cm,y=.36cm,font=\scriptsize]
  \roosterleeg{-3}{1}{3}{-3}{1}{3}
  \draw[accent,thick] (-3,3) -- (0,0) -- (3,3);
  \node[below,text=black] at (0,-3.2) {$|x|$};
\end{tikzpicture}\hspace{3mm}
\begin{tikzpicture}[x=.36cm,y=.36cm,font=\scriptsize]
  \roosterleeg{-3}{1}{3}{-3}{1}{3}
  \kromme{.34}{3}{1/\x} \kromme{-3}{-.34}{1/\x}
  \node[below,text=black] at (0,-3.2) {$\frac1x$};
\end{tikzpicture}}

\begin{theorie}
\renewcommand{\arraystretch}{1.25}
\begin{tabular}{@{}l@{\hspace{8mm}}ccccc@{}}
$f(x)$ & $\formule{x^2}$ & $\formule{x^3}$ & $\formule{\sqrt{x}}$ & $\formule{|x|}$ & $\formule{\frac1x}$ \\
$\dom f$ & $\R$ & $\R$ & $\left[0, +\infty\right[$ & $\R$ & $\R \setminus \{0\}$ \\
$\ber f$ & $\left[0, +\infty\right[$ & $\R$ & $\left[0, +\infty\right[$ & $\left[0, +\infty\right[$ & $\R \setminus \{0\}$ \\
symmetrie & even & oneven & geen & even & oneven
\end{tabular}
\end{theorie}

\begin{oefeningen}
\oef{h2-basis}
\end{oefeningen}

%========================================================================================
\section{Verschuiven}

\begin{theorie}
\begin{stelling}{Verschuiving}
De grafiek van $\formule{g(x) = f(x - c) + d}$ is de grafiek van $f$, verschoven over $c$ naar rechts en $d$ naar boven. Een punt $(x, y)$ gaat naar $\formule{(x + c, y + d)}$.
\end{stelling}
\end{theorie}

Let op het teken: in $f(x - 3)$ schuift de grafiek naar \textit{rechts}. Wat $f$ eerst in $0$ deed, doet $g$ nu in $x = 3$, want $g(3) = f(0)$.

\voorbeeld{}{Schets de grafiek van $g(x) = \sqrt{x + 3} - 1$, en geef het domein en het bereik.\label{vb:h2-verschuiven}}{%
\naast{%
\uitleg{Schrijf $x + 3$ als $x - (-3)$: dat is $c = -3$, en $d = -1$. Volg het beginpunt $(0, 0)$ van $\sqrt{x}$.}
\begin{stappenvrij}{4.3cm}{3.3cm}
c = -3,\ d = -1 & 3 naar links, 1 naar onder. \\
(0, 0) \mapsto (-3, -1) & Het beginpunt. \\
\dom g = \left[-3, +\infty\right[ & Ook 3 naar links. \\
\ber g = \left[-1, +\infty\right[ & Ook 1 naar onder.
\end{stappenvrij}}{%
\begin{tikzpicture}[x=.5cm,y=.5cm,font=\scriptsize]
  \rooster{-3}{1}{4}{-1}{1}{3}
  \draw[grijs,dashed,thick] plot[domain=0:4,samples=60,smooth] (\x,{sqrt(\x)});
  \kromme{-3}{4}{sqrt(\x+3)-1}
  \vol{(-3,-1)}
\end{tikzpicture}}}

\nujij{Waar ligt de top van de parabool $y = (x - 2)^2 + 1$?}{De top $(0, 0)$ van $y = x^2$ schuift 2 naar rechts en 1 naar boven: $(2, 1)$.}

\begin{oefeningen}
\oef{h2-verschuiven}

\needspace{8\baselineskip}
\oef{h2-verschuiven-grafiek}
\end{oefeningen}

%========================================================================================
\section{Uitrekken en spiegelen}

\begin{theorie}
\begin{stelling}{Uitrekken}
De grafiek van $\formule{g(x) = a \cdot f(x)}$ is die van $f$, \textbf{verticaal} uitgerekt met factor $a$: $\formule{(x, y) \mapsto (x, ay)}$.

De grafiek van $\formule{h(x) = f(bx)}$ is die van $f$, \textbf{horizontaal} uitgerekt met factor $\frac1b$: $\formule{(x, y) \mapsto \left(\frac{x}{b}, y\right)}$.
\end{stelling}
\begin{stelling}{Spiegelen}
De grafiek van $\formule{-f(x)}$ is het spiegelbeeld van die van $f$ om de $x$-as, de grafiek van $\formule{f(-x)}$ het spiegelbeeld om de $y$-as.
\end{stelling}
\end{theorie}

Ook hier werkt een factor binnen de haakjes \textit{omgekeerd}: $f(2x)$ maakt de grafiek twee keer smaller, want $h(1) = f(2)$: wat $f$ in $2$ doet, doet $h$ al in $1$.

\voorbeeld{}{De grafiek van $f$ is de gebroken lijn door $(-2, 0)$, $(0, 2)$, $(2, 0)$ en $(4, -2)$. Teken de grafieken van $g(x) = 2f(x)$ en $h(x) = f(2x)$.\label{vb:h2-rekken}}{%
\uitleg{Transformeer de vier hoekpunten en verbind ze. Bij $g$ verdubbelen de $y$-waarden, bij $h$ halveren de $x$-waarden.}
\begin{stappenbreed}
g\colon\ (-2, 0),\ (0, 4),\ (2, 0),\ (4, -4) & $(x, y) \mapsto (x, 2y)$: twee keer zo hoog. \\
h\colon\ (-1, 0),\ (0, 2),\ (1, 0),\ (2, -2) & $(x, y) \mapsto (\frac{x}{2}, y)$: twee keer zo smal.
\end{stappenbreed}
\figuren{%
\begin{tikzpicture}[x=.42cm,y=.42cm,font=\scriptsize]
  \rooster{-2}{1}{4}{-4}{1}{4}
  \draw[grijs,dashed,thick] (-2,0) -- (0,2) -- (2,0) -- (4,-2);
  \draw[accent,thick] (-2,0) -- (0,4) -- (2,0) -- (4,-4);
  \node[right,text=accent] at (.4,3.6) {$g$};
\end{tikzpicture}\hspace{1cm}
\begin{tikzpicture}[x=.42cm,y=.42cm,font=\scriptsize]
  \rooster{-2}{1}{4}{-4}{1}{4}
  \draw[grijs,dashed,thick] (-2,0) -- (0,2) -- (2,0) -- (4,-2);
  \draw[accent,thick] (-1,0) -- (0,2) -- (1,0) -- (2,-2);
  \node[right,text=accent] at (1.6,-2.3) {$h$};
\end{tikzpicture}}}

\nujij{Welk domein en bereik heeft $g(x) = \sqrt{-x}$? En $k(x) = -\sqrt{x}$?}{$g$ is $\sqrt{x}$ gespiegeld om de $y$-as: $\dom g = \left]-\infty, 0\right]$, $\ber g = \left[0, +\infty\right[$. $k$ is gespiegeld om de $x$-as: $\dom k = \left[0, +\infty\right[$, $\ber k = \left]-\infty, 0\right]$.}

\begin{oefeningen}
\oef{h2-rekken}

\oef{h2-spiegelen}
\end{oefeningen}

%========================================================================================
\section{Transformaties combineren}

Bij een voorschrift als $g(x) = a \cdot f\bigl(b(x - c)\bigr) + d$ voer je de transformaties na elkaar uit. Werk van binnen naar buiten: eerst wat bij $x$ gebeurt (horizontaal), dan wat met $f(\ldots)$ gebeurt (verticaal).

\voorbeeld{}{Schets $g(x) = -2\sqrt{x + 1} + 3$ vanuit $\sqrt{x}$, en geef het domein en het bereik.\label{vb:h2-combineren}}{%
\naast{%
\uitleg{Volg het beginpunt $(0, 0)$ door alle stappen.}
\begin{stappenvrij}{3.9cm}{3.7cm}
\sqrt{x + 1} & 1 naar links: $(-1, 0)$. \\
2\sqrt{x + 1} & Verticaal $\times 2$: $(-1, 0)$. \\
-2\sqrt{x + 1} & Spiegelen om de $x$-as. \\
-2\sqrt{x + 1} + 3 & 3 naar boven: $(-1, 3)$. \\
\dom g = \left[-1, +\infty\right[ & \\
\ber g = \left]-\infty, 3\right] & De grafiek daalt vanaf $3$.
\end{stappenvrij}}{%
\begin{tikzpicture}[x=.5cm,y=.5cm,font=\scriptsize]
  \rooster{-1}{1}{5}{-2}{1}{3}
  \draw[grijs,dashed,thick] plot[domain=0:5,samples=60,smooth] (\x,{sqrt(\x)});
  \kromme{-1}{5}{-2*sqrt(\x+1)+3}
  \vol{(-1,3)}
\end{tikzpicture}}}

Staat er een factor bij $x$, zet die dan eerst buiten haakjes: $\sqrt{2x - 6} = \sqrt{2(x - 3)}$ is $\sqrt{x}$ horizontaal $\times \frac12$ en daarna 3 (niet 6) naar rechts.

\begin{werkwijze}
Schrijf het voorschrift als $a \cdot f\bigl(b(x - c)\bigr) + d$. Pas toe in deze volgorde: horizontaal uitrekken met $\frac1b$ (en spiegelen als $b < 0$), verschuiven over $c$ naar rechts, verticaal uitrekken met $a$ (en spiegelen als $a < 0$), verschuiven over $d$ naar boven. Volg een bijzonder punt, en controleer met één functiewaarde.
\end{werkwijze}

\begin{oefeningen}
\oef{h2-combineren-schets}
\end{oefeningen}

Omgekeerd kun je uit een grafiek het voorschrift halen, als je weet van welke basisfunctie ze komt.

\voorbeeld{}{De parabool hiernaast heeft top $(-1, 2)$ en gaat door $(1, 0)$. Geef het voorschrift.\label{vb:h2-voorschrift}}{%
\naast{%
\uitleg{Een parabool is $x^2$, verschoven naar de top en verticaal uitgerekt: $y = a(x - c)^2 + d$. Bepaal $a$ met het tweede punt.}
\begin{stappenvrij}{4cm}{3.6cm}
y = a(x + 1)^2 + 2 & Top $(-1, 2)$. \\
0 = a \cdot 2^2 + 2 & Punt $(1, 0)$ invullen. \\
a = -\tfrac12 & \\
y = -\tfrac12(x + 1)^2 + 2 &
\end{stappenvrij}}{%
\begin{tikzpicture}[x=.5cm,y=.5cm,font=\scriptsize]
  \rooster{-4}{1}{2}{-1}{1}{3}
  \kromme{-3.5}{1.5}{-(\x+1)^2/2+2}
  \vol{(-1,2)} \vol{(1,0)}
\end{tikzpicture}}}

\begin{oefeningen}
\oef{h2-combineren}

\oef{h2-voorschrift-kort}
\end{oefeningen}

%========================================================================================
\opdrachtenkop

\begin{oefeningen}
\needspace{10\baselineskip}
\oef{h2-o1}

\needspace{9\baselineskip}
\oef{h2-o10}

\oef{h2-o2}

\needspace{9\baselineskip}
\oef{h2-o11}

\needspace{7\baselineskip}
\oef{h2-o14}

\needspace{6\baselineskip}
\oef{h2-o12}

\needspace{8\baselineskip}
\oef{h2-o13}

\needspace{10\baselineskip}
\oef{h2-o16}

\needspace{9\baselineskip}
\oef{h2-o3}

\needspace{7\baselineskip}
\oef{h2-o20}

\needspace{8\baselineskip}
\oef{h2-o4}

\needspace{10\baselineskip}
\oef{h2-o6}

\needspace{9\baselineskip}
\oef{h2-o17}

\needspace{11\baselineskip}
\oef{h2-o5}

\needspace{9\baselineskip}
\oef{h2-o21}

\needspace{8\baselineskip}
\oef{h2-o24}

\needspace{9\baselineskip}
\oef{h2-o18}

\needspace{10\baselineskip}
\oef{h2-o25}

\needspace{10\baselineskip}
\oef{h2-o7}

\needspace{10\baselineskip}
\oef{h2-o27}

\needspace{8\baselineskip}
\oef{h2-o26}

\oef{h2-o8}

\needspace{8\baselineskip}
\oef{h2-o28}

\oef{h2-o9}
\end{oefeningen}

%----------------------------------------------------------------------------------------
% KOPPEL: verbanden met andere leerstof
%----------------------------------------------------------------------------------------
\needspace{20\baselineskip}
\opdrachtsectie{Koppel}

\lkr{Goniometrische functies, Exponentiële en logaritmische functies en Statistiek: kijk na wat al gezien is. Het pakket legt de volgorde niet vast.}

\begin{oefeningen}
\oef{h2-koppel-1}

\oef{h2-koppel-2}

\oef{h2-koppel-3}
\end{oefeningen}

%----------------------------------------------------------------------------------------
\opdrachtsectie{Test jezelf}

\begin{oefeningen}
\oef{h2-t5}

\oef{h2-t1}

\oef{h2-t2}

\oef{h2-t6}

\oef{h2-t3}

\oef{h2-t7}

\oef{h2-t4}
\end{oefeningen}

\herhaalschema{%
  \fzv{oef:h2-t5}{oef:h2-basis}%
  \fv{oef:h2-t1}{vb:h2-verschuiven}{oef:h2-verschuiven}%
  \fv{oef:h2-t2}{vb:h2-rekken}{oef:h2-rekken}%
  \fv{oef:h2-t6}{vb:h2-rekken}{oef:h2-o10}%
  \fv{oef:h2-t3}{vb:h2-combineren}{oef:h2-combineren}%
  \fv{oef:h2-t7}{vb:h2-combineren}{oef:h2-combineren-schets}%
  \fv{oef:h2-t4}{vb:h2-voorschrift}{oef:h2-voorschrift-kort}}

\ifzelfstudie
\section*{Maak je eigen samenvatting}
Sluit dit boek en schrijf in je schrift: de grafieken van de vijf basisfuncties met domein en bereik, wat $f(x - c) + d$, $a f(x)$, $f(bx)$, $-f(x)$ en $f(-x)$ met een grafiek doen, en in welke volgorde je combineert. Vergelijk daarna met de kaders.
\fi
